\documentclass[11pt,a4paper]{article}

% ── Packages ────────────────────────────────────────────────────────────────
\usepackage[utf8]{inputenc}
\usepackage[T1]{fontenc}
\usepackage[margin=1in]{geometry}
\usepackage{amsmath, amssymb, amsthm, mathtools, bm}
\usepackage{graphicx}
\usepackage{booktabs}
\usepackage{tabularx}
\usepackage{enumitem}
\usepackage{hyperref}
\usepackage{xcolor}
\usepackage{listings}
\usepackage{float}

% ── Listings setup (for code-style snippets) ────────────────────────────────
\lstset{
  basicstyle=\ttfamily\footnotesize,
  breaklines=true,
  columns=fullflexible,
  keepspaces=true,
  frame=single,
  showstringspaces=false,
  numberstyle=\tiny\color{gray},
  commentstyle=\color{gray!70!black}\itshape,
}

% ── Hyperref colours ────────────────────────────────────────────────────────
\hypersetup{
  colorlinks=true,
  linkcolor=blue!50!black,
  citecolor=blue!50!black,
  urlcolor=blue!60!black,
}

% ── Convenience macros ──────────────────────────────────────────────────────
\newcommand{\Utarget}{U_{\text{target}}}
\newcommand{\Ucirc}{U_{\text{circ}}}
\newcommand{\BOS}{\ifmmode\langle\text{\textsc{bos}}\rangle\else$\langle$\textsc{bos}$\rangle$\fi}
\newcommand{\STOP}{\ifmmode\langle\text{\textsc{stop}}\rangle\else$\langle$\textsc{stop}$\rangle$\fi}
\newcommand{\Real}{\mathbb{R}}
\newcommand{\Complex}{\mathbb{C}}
\newcommand{\Expect}{\mathbb{E}}
\newcommand{\softmax}{\operatorname{softmax}}
\newcommand{\KL}{\mathrm{KL}}
\newcommand{\E}{\mathrm{e}}

\title{\textbf{Hybrid-Action Generative Quantum Eigensolver:\\
A Transformer-Based Reinforcement Learning Pipeline\\
for Quantum Circuit Compilation}}
\author{Ayman Tarig}

\begin{document}
\maketitle

% ════════════════════════════════════════════════════════════════════════════
\section{Introduction and Problem Definition}
% ════════════════════════════════════════════════════════════════════════════

\subsection{Quantum compilation}

Let $\mathcal{H} = (\Complex^2)^{\otimes n}$ denote the $n$-qubit Hilbert space of dimension $d = 2^n$, and let $\Utarget \in U(d)$ be a target unitary that we wish to implement on a near-term device. Such a device admits only a finite native gate set; we adopt the hardware-aware basis
\begin{equation}
\mathcal{G}
=
\Bigl\{
\mathrm{R}_x(\theta)_q,\;
\mathrm{R}_y(\theta)_q,\;
\mathrm{R}_z(\theta)_q,\;
\mathrm{SX}_q,\;
\mathrm{CNOT}_{c,t}
\;\Big|\;
q \in [n],\; (c,t) \in [n]^2,\; c \neq t
\Bigr\},
\end{equation}
in which the single-qubit rotations are parametric while SX and CNOT are fixed. Compilation amounts to finding an ordered sequence $\bm{g} = (g_1, \dots, g_T)$ with $g_i \in \mathcal{G}$ and a vector of rotation angles $\bm{\theta} \in \Real^T$ whose composition
\begin{equation}
\Ucirc(\bm{g}, \bm{\theta})
\;=\;
g_T(\theta_T)\, g_{T-1}(\theta_{T-1})\, \cdots\, g_1(\theta_1)
\end{equation}
approximates $\Utarget$ as measured by the \emph{process fidelity}
\begin{equation}
\boxed{\;\;
F(\Utarget, \Ucirc)
\;=\;
\frac{\bigl|\,\mathrm{tr}(\Utarget^{\dagger}\,\Ucirc)\,\bigr|^{2}}{d^{2}}
\;\;\in [0,1].
\;\;}
\label{eq:process-fidelity}
\end{equation}

\subsection{Why a multi-objective formulation?}

Optimising fidelity alone is degenerate: a sufficiently long random sequence of pool gates can drive $F \to 1$, but the resulting circuit is unusable on real hardware. We therefore optimise three competing quantities at once:
\begin{itemize}[leftmargin=2em]
  \item maximise the process fidelity $F \in [0, 1]$;
  \item minimise the circuit \emph{depth} $D$, defined as the longest qubit-wise dependency chain;
  \item minimise the \emph{CNOT count} $C$, which is the dominant noise contributor.
\end{itemize}
For each fidelity threshold, the Pareto-optimal set in $(F, D, C)$ space contains the shallowest and lowest-CNOT solutions. The trained policy is therefore evaluated not on a single best circuit but on the entire non-dominated front it discovers.

\subsection{Why reinforcement learning, and why hybrid actions?}

The combinatorial choice of \emph{which gate at which qubit} makes generic gradient-based variational compilation brittle: one must commit to a fixed gate template (an \emph{ansatz}) and optimise only its angles. Reinforcement learning sidesteps the ansatz by treating circuit synthesis as sequence generation under a learned, sampling-based policy.

Earlier RL approaches discretise each parametric gate into a finite family $\{\mathrm{R}_z(\pi/8), \mathrm{R}_z(\pi/4), \dots\}$, which multiplies the vocabulary by the angle resolution and forces a post-hoc continuous tuning step. We instead let the policy emit \emph{both} a token and a real-valued angle at every step --- a \textbf{hybrid (discrete + continuous) action} --- thereby unifying structural and parametric search under one objective.

% ════════════════════════════════════════════════════════════════════════════
\section{Detailed Methodology}
% ════════════════════════════════════════════════════════════════════════════

\subsection{Vocabulary and action space}
\label{sec:vocab}

The \emph{operator pool} $\mathcal{P}$ is enumerated once at start-up. For each qubit $q \in \{0, \dots, n-1\}$ the pool contains, in order, the configured rotation gates (the default configuration uses $\mathrm{R}_z$ and $\mathrm{R}_y$ on every qubit) followed by an SX gate, and finally one CNOT for every ordered pair $(c, t)$ with $c \neq t$. Two bookkeeping tokens are prepended to form the vocabulary
\begin{equation}
\mathcal{V} \;=\;
\{\,
\BOS = 0,\;
\STOP = 1,\;
g_2,\; g_3,\; \dots,\; g_{V-1}
\,\},
\qquad V = 2 + |\mathcal{P}|.
\end{equation}
At every sequence position $t \in \{1,\dots, T_{\max}\}$ the policy emits a joint action
\begin{equation}
a_t = (a^d_t,\, a^c_t),
\qquad
a^d_t \in \mathcal{V},\quad a^c_t \in \Real,
\end{equation}
where the continuous component $a^c_t$ is interpreted as a rotation angle if and only if $a^d_t$ is one of $\mathrm{R}_x, \mathrm{R}_y, \mathrm{R}_z$, and is otherwise discarded.

\subsection{Transformer backbone}
\label{sec:transformer}

The backbone is a standard decoder-only GPT-2 with tied input/output token embeddings. Four sizes are exposed:

\begin{center}
\small
\begin{tabular}{lccc}
\toprule
size & layers ($L$) & heads ($H$) & embed dim ($d_{\!e}$)\\
\midrule
\texttt{tiny}   & 2  & 2  & 128 \\
\texttt{small}  & 6  & 6  & 384 \\
\texttt{medium} & 12 & 12 & 768 \\
\texttt{large}  & 24 & 16 & 1024\\
\bottomrule
\end{tabular}
\end{center}

\paragraph{Forward pass.}
Given a token sequence $x = (x_0,\dots,x_{T-1}) \in \mathcal{V}^T$,
\begin{equation}
h_0 \;=\; W_e[x] + W_p[0..T-1]
\quad\xrightarrow{\text{Block} \times L}\quad
h_L \;=\; \mathrm{LayerNorm}(\,\cdot\,),
\end{equation}
where each block is the canonical pre-norm transformer block
\begin{equation}
\tilde{h} \;=\; h + \mathrm{MHA}\bigl(\mathrm{LN}(h),\,\text{causal mask}\bigr),
\qquad
h' \;=\; \tilde{h} + \mathrm{MLP}\bigl(\mathrm{LN}(\tilde{h})\bigr).
\end{equation}
The MLP uses the GPT-2 ``new GeLU'' activation
$\mathrm{GeLU}(z) = \tfrac{1}{2} z\,(1 + \tanh(\sqrt{2/\pi}(z + 0.044715 z^3)))$
and a $4\times$ widening factor.

\paragraph{Two heads.}
On the final hidden state $h_L \in \Real^{B\times T\times d_{\!e}}$ we attach two heads:
\begin{align}
\textbf{Discrete head:}\quad
\bm{\ell}_{t} &\;=\; W_e^{\top}\,h_{L,t}
\;\in\;\Real^{V}, \quad\text{(weights \emph{tied} with the input embedding),}\\
\textbf{Continuous head:}\quad
\bigl(\mu_t,\, \log\sigma_t\bigr)
&\;=\; W_{\mathrm{co}}\;\mathrm{GeLU}(W_{\mathrm{cf}}\,h_{L,t})
\;\in\;\Real^{2}.
\end{align}
The log-standard-deviation is clipped to $[-5, 2]$ to prevent two pathologies: collapse to a delta-policy ($\sigma\to 0$), and the early-training log-prob explosion that occurs when $\sigma$ is allowed to grow unboundedly.

\paragraph{Sampling.}
The discrete sample is drawn from a temperature-scaled softmax,
\begin{equation}
\Pr(a^d_t = v\mid s_t)
\;=\;
\softmax\bigl(-\beta\,\bm{\ell}_{t}\bigr)_v,
\label{eq:beta-softmax}
\end{equation}
where $\beta$ is the \emph{inverse temperature} (Section~\ref{sec:scheduler}). The continuous sample is unbounded:
\begin{equation}
a^c_t \;\sim\; \mathcal{N}\bigl(\mu_t,\,\E^{2\log\sigma_t}\bigr).
\end{equation}
Because rotation angles are $2\pi$-periodic, no squashing transform (e.g.\ Tanh) is needed; using a plain Gaussian also keeps the log-density exact (no Tanh-correction term) and the PPO importance ratio numerically stable.


\subsection{GRPO-style PPO with hybrid actions}
\label{sec:grpo}

\paragraph{Joint per-position log-probability.}
Let $m_t \in \{0,1\}$ indicate whether token $a^d_t$ is parametric (i.e.\ an $\mathrm{R}_x$, $\mathrm{R}_y$, or $\mathrm{R}_z$). The joint log-probability of action $a_t$ is then
\begin{equation}
\log \pi_{\bm{\phi}}(a_t \mid s_t)
\;=\;
\underbrace{\log\softmax(-\beta\,\bm{\ell}_t)_{a^d_t}}_{\text{discrete}}
\;+\;
m_t\,\underbrace{\log
\mathcal{N}\bigl(a^c_t \mid \mu_t,\, \sigma_t^2\bigr)}_{\text{continuous}}.
\label{eq:joint-logp}
\end{equation}

\paragraph{Sequence-level objective.}
Per-position log-probs are averaged over the \emph{valid} prefix --- positions strictly before the first \STOP\ --- so that comparisons across variable-length samples are well defined:
\begin{equation}
\log \pi_{\bm{\phi}}(\bm{a} \mid \bm{s})
\;=\;
\frac{1}{L}\sum_{t=1}^{L}
\log \pi_{\bm{\phi}}(a_t \mid s_t),
\qquad
L = \min\{t : a^d_t = \STOP\}.
\end{equation}

\paragraph{Group-relative advantages (GRPO).}
For each rollout group --- one batch of $B$ samples drawn from the same policy snapshot --- we normalise the cost relative to the group's mean and standard deviation:
\begin{equation}
A_i
\;=\;
\frac{\bar{c} - c_i}{\mathrm{std}(c) + \varepsilon},
\qquad
\bar{c} = \tfrac{1}{B}\sum_i c_i, \quad \varepsilon = 10^{-8}.
\label{eq:grpo-adv}
\end{equation}
This is the GRPO baseline. It dispenses with a learned critic by exploiting the fact that all samples in a rollout share the same starting state. Lower-cost samples receive positive advantage and higher-cost samples receive negative advantage; subtracting the group mean keeps the policy gradient unbiased, while dividing by the standard deviation makes the loss scale-invariant across reward shapes.

\paragraph{Clipped surrogate.}
With $r_i = \exp\bigl(\log \pi_{\bm{\phi}}(\bm{a}_i) - \log \pi_{\bm{\phi}_{\text{old}}}(\bm{a}_i)\bigr)$ the importance ratio against the behaviour policy, the PPO loss is
\begin{equation}
\mathcal{L}_{\text{PPO}}
\;=\;
-\,\Expect_i
\Bigl[
\min\bigl(
r_i\,A_i,\;
\mathrm{clip}(r_i,\, 1-\epsilon,\, 1+\epsilon)\,A_i
\bigr)
\Bigr],
\qquad \epsilon = 0.2.
\label{eq:ppo-clipped}
\end{equation}

\paragraph{Entropy regularisation.}
Two entropy terms are added with separately tunable weights $\alpha_d, \alpha_c$ (both default to $5\!\cdot\!10^{-3}$):
\begin{equation}
\mathcal{L}_{\text{tot}}
\;=\;
\mathcal{L}_{\text{PPO}}
\;-\;
\alpha_d\,\bar{H}_{\text{cat}}
\;-\;
\alpha_c\,\bar{H}_{\text{Gauss}},
\end{equation}
where, with $\mathcal{T}$ denoting the set of valid (non-padding) positions,
\begin{align}
\bar{H}_{\text{cat}}
&= \Expect\Bigl[
-\textstyle\sum_{v} p_{t,v}\log p_{t,v},\quad
p = \softmax(-\beta\,\bm{\ell}_t)
\Bigr]_{t \in \mathcal{T}},\\
\bar{H}_{\text{Gauss}}
&= \Expect\Bigl[
\log\sigma_t + \tfrac{1}{2}\log(2\pi\E)
\Bigr]_{t \in \mathcal{T}}.
\end{align}
The categorical entropy uses a careful \texttt{where(p > 0, log\_p, 0)} guard \emph{before} multiplying by $p$, since action-mask slots have $p=0$ and JAX autograd evaluates both branches of \texttt{where} symbolically.

\paragraph{Optimisation.}
The full policy parameters $\bm{\phi}$ are updated by AdamW with $\eta = 10^{-4}$, weight decay $0.01$, and global gradient-norm clipping at $1.0$. PPO updates run for \texttt{steps\_per\_epoch} (default $4$) mini-batch passes over the replay buffer per training epoch.

\subsection{Reward function design}
\label{sec:reward}

The agent's per-sample \emph{cost} (PPO minimises cost) is $c_i = -R_i$.

\subsubsection{QASER formula}

The reward, inspired by \texttt{arXiv:2511.16272}, is
\begin{equation}
R^{\text{qaser}}
\;=\;
\Biggl(
w_D\,\frac{M_D}{D + 1} +
w_C\,\frac{M_C}{C + 1} +
w_G\,\frac{M_G}{G + 1}
\Biggr)^{F}
\;-\; 1,
\label{eq:reward-qaser}
\end{equation}
where $G$ is the total gate count and $M_D, M_C, M_G$ are running maxima of depth, CNOT count, and total gate count, initialised from the worst case observed in the first rollout (or from configured initial values) and updated thereafter. The per-axis weights are tunable; the configured defaults are $w_D = 1.0$, $w_C = 2.0$, $w_G = 1.5$, with $w_C$ the largest because two-qubit gates dominate hardware noise. Two design choices distinguish this from a flat structural penalty:
\begin{itemize}[leftmargin=2em]
  \item \textbf{Multiplicative base, $F$ as exponent.} For low fidelity the base raised to a small $F$ collapses toward $1$, so after subtracting $1$ the reward $\to 0$ regardless of structure. The base only begins to matter as $F \to 1$, so structural compactness is rewarded \emph{only at high fidelity} --- precisely the regime in which we want the agent to discriminate.
  \item \textbf{Three structural axes.} Including the total gate count $G$ (the original paper's ``one-qubit-gate weight'') prevents the agent from sprinkling free single-qubit rotations between CNOTs once the fidelity is high; without it, single-qubit gates would carry zero marginal cost.
\end{itemize}
Subtracting $1$ centres the reward so that very-low-fidelity samples do not produce a flat constant that would zero out the GRPO advantage.

\subsubsection{Log-infidelity multiplier}

A configurable extension multiplies the QASER reward by a log-infidelity factor,
\begin{equation}
R \;\leftarrow\; R\,\bigl(1 - \log(1 - F + \varepsilon_{\text{lf}})\bigr),
\label{eq:reward-loginf}
\end{equation}
with a small floor $\varepsilon_{\text{lf}}$ (default $10^{-3}$, set to $0$ to recover the paper-faithful QASER form). At low and mid fidelity the multiplier is $\approx 1$ and behaviour is unchanged, but as $F \to 1$ it grows without bound, providing an unbounded marginal pull that escapes the $F \approx 0.98$ plateau observed on harder targets where the pure $\text{base}^F$ form flattens out. Additionally, a \texttt{fidelity\_floor} (default $0.5$) gates archive insertion: candidates below the floor are not eligible for the Pareto archive, preventing trivially short low-fidelity circuits from being reported as ``non-dominated on compactness''.

\subsection{Continuous parameter optimisation}
\label{sec:angle-opt}

The policy outputs a \emph{distribution} over angles, but for any fixed discrete sequence $\bm{g}$ the optimal angles are recovered by classical gradient descent on the differentiable cost $\theta \mapsto 1 - F\bigl(\Utarget,\,\Ucirc(\bm{g}, \bm{\theta})\bigr)$. We exploit this in two distinct refinement stages.

\subsubsection{Inner refinement (HyRLQAS-style)}

After sampling $(\bm{g}_i, \bm{\theta}_i^{(0)})$ from the policy at every rollout, the angles are polished by a few iterations of L-BFGS, initialised at $\bm{\theta}_i^{(0)}$ at learning rate $0.1$. The default \texttt{inner\_refine\_steps} is $12$ (5--10 is the typical sweet spot for small/medium models, while a tiny model benefits from 12--16 to compensate for limited capacity):
\begin{equation}
\bm{\theta}_i^{\star}
\;=\;
\arg\min_{\bm{\theta}}\;
\bigl[\,1 - F\bigl(\Utarget,\,\Ucirc(\bm{g}_i, \bm{\theta})\bigr)\,\bigr],
\qquad
\bm{\theta}^{(0)}_{\text{init}} = \bm{\theta}^{(0)}_i.
\end{equation}
The reward (Equation~\ref{eq:reward-qaser}) is computed using $F^{\star}_i = F(\Utarget, \Ucirc(\bm{g}_i, \bm{\theta}_i^{\star}))$, so the agent is rewarded not for guessing the optimal angle outright but for proposing an \emph{initial} angle that converges quickly under classical optimisation. This is the HyRLQAS hybrid scheme: the RL policy and the classical optimiser cooperate on different aspects of the same task --- structure versus angles.

A subtlety arises in the bookkeeping. The PPO replay buffer stores the \emph{initial} (RL-sampled) angles $\bm{\theta}_i^{(0)}$, not the refined values $\bm{\theta}_i^{\star}$, so the importance ratio in Equation~\ref{eq:ppo-clipped} is computed against the action that was actually sampled (otherwise the policy gradient would be biased). The Pareto archive, by contrast, stores the refined values, since those are the angles that define the reported circuit. To guard against rare L-BFGS overshoots, refinement is clipped to be monotone non-decreasing in fidelity:
\[
F_i^{\text{report}} = \max\bigl(F_i^{\star},\, F_i^{\text{raw}}\bigr).
\]

\subsubsection{Post-training refinement}

Once RL converges, the entire Pareto archive is re-optimised with a larger budget (\texttt{refinement.steps} $= 100$ L-BFGS steps, \texttt{lr} $= 0.1$). The pipeline is:
\begin{enumerate}[leftmargin=2em]
  \item For each archive entry $p$, optionally run a structural simplifier that
    \begin{itemize}
      \item merges consecutive same-axis rotations on the same qubit, $R_z(a)\,R_z(b) \to R_z(a+b)$;
      \item cancels adjacent identical CNOTs;
      \item commutes-then-merges/cancels through any run of gates that commute with the anchor on its support.
    \end{itemize}
    Merged or cancelled positions are overwritten with \STOP\ so they decode to NOOP downstream. Because the angles are re-optimised afterwards, the simplifier need not bookkeep angle merges precisely.
  \item Re-run L-BFGS on the simplified token sequence with multi-restart: one trajectory is initialised from the inner-refined angles, and the remaining \texttt{num\_restarts} $- 1$ trajectories (default total $16$) are seeded from independent $\mathrm{Uniform}[-\pi, \pi]$ vectors. The best-fidelity restart per circuit is kept. This escapes the highly non-convex parameterised-circuit loss landscape that a single L-BFGS trajectory routinely gets trapped in.
  \item Insert the polished entry into a fresh archive; any entries that become dominated after refinement are evicted.
\end{enumerate}

\subsubsection{The fidelity kernel}

Both refiners share a single JIT-compiled forward pass. The circuit unitary is built by a \texttt{jax.lax.scan} over $T_{\max}$ steps; at each step the running $d \times d$ unitary is updated by either a single-qubit or two-qubit gate \emph{without} ever materialising the full $d\!\times\!d$ embedded operator:
\begin{itemize}[leftmargin=2em]
  \item \textbf{Qubit-local apply.} Applying a $2\times 2$ gate $G$ to qubit $k$ of $U \in \Complex^{d\times d}$ is implemented as a row-permutation gather, two block linear combinations $Y_0 = G_{00}U_0 + G_{01}U_1$ and $Y_1 = G_{10}U_0 + G_{11}U_1$, and an inverse gather. The cost is $O(d^2)$ per gate, versus $O(d^3)$ for a generic $d\!\times\!d$ matmul.
  \item \textbf{CNOT apply.} The same trick is applied on a four-block partitioning by the (control, target) bit pair.
  \item \textbf{Custom VJP for permutation.} Native indexing \texttt{U[perm]} backward-passes via scatter-add; we register a manual VJP that uses a single inverse gather, yielding roughly a 50\% backward-pass speed-up.
\end{itemize}
The forward returns $F = \bigl|\sum_{ij} \overline{(\Utarget)}_{ij}\,(\Ucirc)_{ij}\bigr|^2 / d^2$, an $O(d^2)$ trace via element-wise product that avoids the $O(d^3)$ matmul-then-trace. Single-precision complex (\texttt{c64}) is used on GPU and double precision (\texttt{c128}) on CPU. The same function is wrapped in \texttt{jax.value\_and\_grad} for both inner and post-training refinement.

\subsection{Pareto archive}
\label{sec:pareto}

The archive maintains a Pareto-optimal set on $(F\!\uparrow, D\!\downarrow, C\!\downarrow)$. Dominance is checked with a fidelity tolerance $\delta_F = 10^{-4}$:
\begin{equation}
A \succ B
\;\Longleftrightarrow\;
\bigl(F_A + \delta_F \ge F_B,\, D_A \le D_B,\, C_A \le C_B\bigr)
\;\wedge\;
\bigl(F_A > F_B + \delta_F \;\vee\; D_A < D_B \;\vee\; C_A < C_B\bigr).
\end{equation}
The tolerance prevents float-32 jitter from polluting the front with near-duplicates. Updates use a single broadcasted $O((B+A)^2)$ pass to filter both incoming candidates and surviving archive members in one go.

When the archive exceeds \texttt{max\_archive\_size} $= 500$, the entry with the smallest \emph{crowding distance} is evicted. The crowding distance sums, over each objective, the normalised gap between an entry's left and right neighbours:
\begin{equation}
\mathrm{cd}(p)
\;=\;
\sum_{o \in \{F,D,C\}}
\frac{o(p_{\text{next}}) - o(p_{\text{prev}})}{o_{\max} - o_{\min}}.
\end{equation}
Boundary points receive $\mathrm{cd} = +\infty$ and are never pruned. As a quality summary we report the 2-D hypervolume in $(F, C)$,
\begin{equation}
\mathrm{HV}_{2D}
\;=\;
\sum_{i=1}^{|F^{\text{2D}}|}
\bigl(F_i - F_{i-1}\bigr)\,\bigl(C_{\text{ref}} - C_i\bigr),
\end{equation}
swept along the 2-D front sorted by ascending $C$ and monotone-increasing $F$.

\subsection{Temperature schedule}
\label{sec:scheduler}

The inverse temperature $\beta$ in Equation~\ref{eq:beta-softmax} controls the sharpness of the discrete distribution at sampling time: low $\beta$ yields uniform exploration, high $\beta$ yields near-greedy behaviour. Three schedules are exposed:
\begin{itemize}[leftmargin=2em]
  \item \textbf{Fixed:} $\beta_t = \beta_0$.
  \item \textbf{Linear:} $\beta_{t+1} = \mathrm{clip}(\beta_t + \Delta,\, \beta_{\min}, \beta_{\max})$.
  \item \textbf{Cosine:} $\beta_t = \tfrac{1}{2}(\beta_{\max}+\beta_{\min}) - \tfrac{1}{2}(\beta_{\max}-\beta_{\min})\cos(2\pi t/\Pi)$, with period $\Pi = \max(1, \texttt{max\_epochs}/2)$.
\end{itemize}
The cosine schedule is the default in \texttt{config.yml}: by cycling between exploration and exploitation, it helps the agent escape the early greedy-then-stuck regimes that single-direction annealing tends to produce.

% ════════════════════════════════════════════════════════════════════════════
\section{The Training Pipeline}
% ════════════════════════════════════════════════════════════════════════════

\subsection{Initialisation}

At start-up the trainer performs the following steps:
\begin{enumerate}[leftmargin=2em]
  \item Build the operator pool $\mathcal{P}$ and the vocabulary $\mathcal{V}$ (Section~\ref{sec:vocab}).
  \item Build the target unitary $\Utarget$ from the configured target type (random, random-reachable, Haar-random, brickwork, or loaded from disk) and verify unitarity, $\|\Utarget\Utarget^{\dagger} - I\|_{\infty} < 10^{-10}$.
  \item Run a Qiskit \texttt{transpile} pass on $\Utarget$ at \texttt{optimization\_level=3} to obtain a baseline reference (depth, total gate count, two-qubit gate count).
  \item Initialise the GPT-2 weights from a Gaussian with $\sigma = 0.02$ and construct the AdamW optimiser with global gradient-norm clipping.
  \item Trace and JIT-compile three hot paths:
    \begin{itemize}
      \item the KV-cached rollout used at sampling time;
      \item the PPO update step;
      \item the structural-metrics function (depth, CNOT count, total gate count).
    \end{itemize}
  \item Allocate the replay buffer (FIFO, default size $B_{\max} = 2048$) and the Pareto archive.
\end{enumerate}

\subsection{Per-epoch loop}
\label{sec:per-epoch}

The per-epoch loop is summarised in Algorithm~\ref{alg:epoch}.

\begin{figure}[H]
\centering
\fbox{\begin{minipage}{0.95\linewidth}
\small
\textbf{Algorithm 1 --- One training epoch.}
\begin{enumerate}[leftmargin=2em, label=\textbf{\arabic*.}]
  \item $\beta \gets \texttt{Scheduler.get()}$
  \item $\bm{\tau} \gets \texttt{KVRollout}(\bm{\phi},\, \beta,\, B)$\\
        \hphantom{xx}\textit{// $B$ samples, each $(\bm{x}_i, \bm{\theta}^{(0)}_i, \log p^d_i, \log p^c_i)$}
  \item $\bm{F}^{\text{raw}} \gets \texttt{Fidelity}(\bm{\tau})$
  \item \textbf{if} inner refinement enabled \textbf{then}
    \begin{itemize}[leftmargin=1em, label={}]
      \item $(\bm{F}^{\star}, \bm{\theta}^{\star}) \gets \texttt{LBFGS}(\bm{\tau},\,\text{steps}{=}12)$
      \item Clip: $\bm{F} = \max(\bm{F}^{\star}, \bm{F}^{\text{raw}})$, $\bm{\theta}^{\text{report}} = \mathbb{1}[\text{improved}]\cdot \bm{\theta}^{\star} + (1-\mathbb{1}[\text{improved}])\cdot\bm{\theta}^{(0)}$
    \end{itemize}
    \textbf{else}\;
    $\bm{F} \gets \bm{F}^{\text{raw}}$,\;
    $\bm{\theta}^{\text{report}} \gets \bm{\theta}^{(0)}$
  \item $\bm{D}, \bm{G}, \bm{C} \gets \texttt{StructureMetrics}(\bm{\tau})$
  \item $\bm{c} \gets \texttt{Cost}(\bm{F}, \bm{D}, \bm{C}, \bm{G})$ \quad
        \textit{// $c_i = -R_i$, eq.~\eqref{eq:reward-qaser}}
  \item \textbf{for} $i = 1, \dots, B$:\;
        $\texttt{Buffer.push}(\bm{x}_i, \bm{\theta}^{(0)}_i, c_i, \log p^d_i, \log p^c_i)$
  \item $\texttt{ParetoArchive.update\_batch}(\{(F_i, D_i, G_i, C_i, \bm{x}_i, \bm{\theta}^{\text{report}}_i)\})$
  \item $\texttt{Scheduler.update}(1 - \bm{F})$
  \item \textbf{for} $k = 1, \dots, \texttt{steps\_per\_epoch}$:
    \begin{itemize}[leftmargin=1em, label={}]
      \item \textbf{for} minibatch $\mathcal{B}\subset \textsc{shuffled}(\text{Buffer})$:
        \begin{itemize}[leftmargin=2em, label={}]
          \item $\mathcal{L} \gets \texttt{PPOLoss}(\mathcal{B}, \beta)$ \quad
                \textit{// eq.~\eqref{eq:ppo-clipped}}
          \item $\bm{\phi} \gets \bm{\phi} - \eta\,\nabla_{\bm{\phi}}\mathcal{L}$
                \quad \textit{// AdamW, gradient-norm clipped}
        \end{itemize}
    \end{itemize}
  \item $\texttt{Logger.write}(\text{epoch metrics, Pareto stats, time})$
  \item \textbf{if} \texttt{early\_stop} \textbf{and} $F_{\max} \ge 1 - 10^{-8}$ \textbf{then} \textbf{break}
\end{enumerate}
\end{minipage}}
\caption{One training epoch of the hybrid-action GQE pipeline.}
\label{alg:epoch}
\end{figure}

\subsection{Replay buffer}

The replay buffer stores tuples $(\bm{x}_i, \bm{\theta}^{(0)}_i, c_i, \log p^d_i, \log p^c_i)$ in a bounded FIFO. Once per epoch the dataset adapter takes a snapshot of the buffer and yields shuffled mini-batches over \texttt{steps\_per\_epoch} repetitions (default $4$). Each PPO update therefore sees roughly four times the rollout volume per epoch without re-sampling --- a standard PPO data-efficiency trick.

With a FIFO depth of $2048$ and $B = 1024$ rollouts per epoch, the agent's effective replay window is about two epochs of fresh trajectories. Older samples are discarded so that the importance ratio $r_i$ never drifts too far from $1$, which in turn keeps the PPO clip a soft regulariser rather than a crippling one.

\subsection{Post-training stage}

After the per-epoch loop terminates --- either by reaching \texttt{max\_epochs} or by early-stop on $F \ge 1 - 10^{-8}$ --- the trainer performs three closing steps:
\begin{enumerate}[leftmargin=2em]
  \item \textbf{Best-raw refinement.} The single highest-fidelity raw rollout produced during training is L-BFGS-refined for \texttt{refinement.steps} (default $100$) iterations under the same multi-restart scheme used for archive refinement. This step is independent of the Pareto archive, ensuring that runs which fail to populate the archive (e.g.\ when the configured \texttt{fidelity\_floor} is too high) still report a polished circuit.
  \item \textbf{Archive refinement.} Each archive entry's tokens are optionally simplified, then its angles are L-BFGS-refined with multi-restart; the polished entries are inserted into a fresh archive so that any entries that become dominated after refinement are evicted.
  \item \textbf{Reporting.} A single ``best to report'' circuit is selected: the preferred choice is the Pareto best-by-depth among entries with $F \ge \texttt{fidelity\_threshold}$ (default $0.99$), falling back to best-by-fidelity, and finally to the best raw rollout. The reported circuit is rebuilt as a Qiskit \texttt{QuantumCircuit} and its fidelity is independently verified by materialising the gate matrices through Qiskit's \texttt{Operator} API.
\end{enumerate}

A run-artifact JSON is emitted to \texttt{results/run\_\{n\}q\_\{type\}\_\{stamp\}.json}. It contains the target unitary, the full configuration snapshot, all per-epoch metrics, and a slim representation (without per-circuit gate lists) of the Pareto front.



\end{document}
